(1) 将相干态写成 $|\alpha_0\rangle=e^{-|\alpha_0|^2/2}\sum_{n=0}^\infty\alpha_0^n|n\rangle/\sqrt{n!}$。每个数态演化相位是 $e^{-i\omega(n+1/2)t}$，因此 $|\psi(t)\rangle=e^{-i\omega t/2}|\alpha_0e^{-i\omega t}\rangle$，即相干参数沿复平面匀速旋转且模长守恒。
(2) 由 $x=\sqrt{\hbar/(2m\omega)}(a+a^\dagger)$ 与 $p=-i\sqrt{m\hbar\omega/2}(a-a^\dagger)$，有 $\langle x(t)\rangle=\sqrt{2\hbar/(m\omega)}\,\Re(\alpha_0e^{-i\omega t})$、$\langle p(t)\rangle=\sqrt{2m\hbar\omega}\,\Im(\alpha_0e^{-i\omega t})$。它们严格满足 $\dot{\langle x\rangle}=\langle p\rangle/m$、$\dot{\langle p\rangle}=-m\omega^2\langle x\rangle$。
(3) 相干态的数态分布是 $P_n=e^{-|\alpha_0|^2}|\alpha_0|^{2n}/n!$，为平均 $\bar n=|\alpha_0|^2$ 的泊松分布，故 $\operatorname{Var}(n)=\bar n$。能量期望为 $\langle H\rangle=\hbar\omega(\bar n+1/2)$，其中 $1/2$ 是真空零点能而不是可忽略的归一修正。
(4) 相干态在无阻尼简谐动力学中保持 $\Delta x^2=\hbar/(2m\omega)$、$\Delta p^2=m\hbar\omega/2$，所以 $\Delta x\Delta p=\hbar/2$，取 Heisenberg 下限。其波包中心做经典周期运动而形状与最小不确定度不变，但仍具有非零真空涨落，绝不等同于相空间中的点粒子。
